\documentclass[11pt]{article}
\usepackage[utf8]{inputenc}
\usepackage{amsmath,amssymb}
\usepackage[margin=1in]{geometry}
\usepackage[hidelinks]{hyperref}
\usepackage{xurl}
\usepackage{setspace}

\title{Charge Densities in Real Space:\\
RealQM and the Realist Quest Carried by Schr\"odinger}
\author{}
\date{}
% ANONYMIZED for double-anonymous review (Synthese);
% author, affiliation, and identifying reference details are on the separate Title Page.

\begin{document}
\maketitle

\begin{abstract}
The central question of realism in atomic physics is whether the microscopic world is of the same form as
the macroscopic --- extended matter in ordinary three-dimensional space, only of smaller size --- or is
something categorically other. Schr\"odinger's wave mechanics began as a bet on the first answer, reading
$\psi\bar\psi$ as a real charge density in space; the Bohr--Born--Heisenberg--Dirac program settled on the
second, and the realist reading was set aside. RealQM takes up Schr\"odinger's quest. We compare the
ontology of an atomic system of $N$ electrons under two models. Standard quantum
mechanics treats the electrons as \emph{identical particles} and represents the system by a probability
amplitude over their configurations in a $3N$-dimensional configuration space, antisymmetric under
exchange. \emph{RealQM} represents it by $N$ unit charge densities, each carried on its own region
$\Omega_i$ of ordinary three-dimensional space, the regions non-overlapping
($\Omega_i\cap\Omega_j=\varnothing$). Two long-standing metaphysical problems --- the \emph{arena} of the
wave function and the \emph{individuality} of identical particles --- attach to the first model, one to
each of its two features: the $3N$-dimensional domain and the antisymmetry. The second model has neither
feature and so incurs neither problem. It is a primitive ontology in three-space (no configuration-space
object exists), and it \emph{individuates} its electrons absolutely, by the regions they occupy (the
Principle of the Identity of Indiscernibles is satisfied); the two virtues are one fact, the non-overlap of
the $\Omega_i$ doing double duty. We then argue that the decisive asset is computational: the standard
amplitude is intractable beyond very small $N$, forcing every practical method to retreat from the
$3N$-dimensional object toward effective three-dimensional ones (density-functional theory makes the
three-dimensional density fundamental in practice), whereas in RealQM the three-dimensional charge
densities \emph{are} the objects of computation, and the same functional has been used to compute atoms
across the periodic table, molecular bonds and hydrogen bonds, and --- reading the nucleus as protons and
electrons bound by Coulomb alone --- nuclear binding energies. The metaphysics is thus not posited for
interpretive tidiness but carried by a working, empirically tested method. We place the proposal against
the point-particle ontology of the Standard Model, answer the main objections, and bound the claim.
\end{abstract}

\setstretch{1.04}

\section{Introduction: the realist question and Schr\"odinger's quest}
\label{sec:intro}

The central question of realism in atomic physics can be put in one line: \emph{is the microscopic of the
same form as the macroscopic --- extended matter in three-dimensional space, only of smaller size --- or is
it categorically other?} On the first answer an atom is a small physical object of the same kind as a large
one, picturable in ordinary space; on the second it is not a thing in space at all but an element of an
abstract formalism whose only contact with space is through the statistics of measurement.

Wave mechanics began on the first side. In 1926 Schr\"odinger read the wave function of an electron as a
continuous distribution of electric charge in ordinary space, $\rho=e\,\psi\bar\psi$, and sought an
\emph{anschaulich} --- intuitive, visualizable --- account of atomic phenomena in terms of such extended
waves, in deliberate opposition to the unvisualizable discontinuities of the older quantum theory
\cite{Schrodinger1926,deRegt1997}. On this reading an atom is a standing charge-density pattern in
three-space: a picture one can draw, of the same form as a macroscopic charge distribution, only atomic in
size.

Within a year the realist reading was displaced. Born recast $\psi\bar\psi$ as a probability density
\cite{Born1926}; Heisenberg's matrix mechanics and uncertainty relations, Bohr's complementarity, and
Dirac's abstract transformation theory turned the formalism into an instrument for computing measurement
probabilities rather than a picture of what there is. Decisive for our purposes is \emph{why} the
charge-density reading had to go: for $N$ electrons the wave function lives not on ordinary space but on a
$3N$-dimensional configuration space, where it cannot be a charge density in $\mathbb R^3$. Anschaulichkeit
was surrendered at exactly the point where the domain left three-space --- the answer to the realist
question turned on the arena of the wave function.

RealQM \cite{RealQM} takes up the abandoned quest. It is a reformulation of the Schr\"odinger equation for
atomic physics that keeps Schr\"odinger's charge densities in ordinary three-space, but resolves the
many-electron system into $N$ \emph{separate} unit densities on non-overlapping regions, so that the
charge-density realism which failed on $\mathbb R^{3N}$ becomes viable again on $\mathbb R^3$. It is,
potentially, a revolutionary change of formulation: the same atomic, molecular, and nuclear physics,
computed and visualized as extended charge in real space, microphysics of the same form as macrophysics.
This paper sets out its ontology, asks what philosophy gains from it --- two classic problems dissolved
together (\S\S\ref{sec:two}--\ref{sec:ind}) --- and asks what keeps it from being one more homemade picture:
that it computes (\S\ref{sec:comp}).

\section{Two problems in one formula}
\label{sec:two}

An atomic system of $N$ electrons admits two descriptions. Standard quantum mechanics treats the electrons
as \emph{identical particles} and represents the system by one probability amplitude over their joint
configurations, in a space of $3N$ dimensions. RealQM \cite{RealQM} treats them as $N$ \emph{individual
charge densities}, each carried on its own region $\Omega_i\subset\mathbb R^3$, in ordinary
three-dimensional space (\S\ref{sec:po}). This contrast --- identical particles in a high-dimensional
configuration space against individuated densities in ordinary space --- is the whole of what follows. We
take the standard side first.

Write its state:
\begin{equation}
\Psi:\ \mathbb R^{3N}\to\mathbb C,\qquad
\Psi(\dots,\mathbf x_i,\dots,\mathbf x_j,\dots)=-\,\Psi(\dots,\mathbf x_j,\dots,\mathbf x_i,\dots).
\label{eq:std}
\end{equation}
The domain is $\mathbb R^{3N}$; the minus sign is the \emph{antisymmetry} of the amplitude under exchange of
any two particles --- Fermi--Dirac statistics, the Pauli principle in its formal guise. Two metaphysical
difficulties follow directly, one from each of these two features of \eqref{eq:std}.

\emph{The arena}, from the domain $\mathbb R^{3N}$. If $\Psi$ is the fundamental object --- wave-function
realism \cite{Albert1996,NeyAlbert2013,Ney2021} --- then ordinary three-space is not fundamental but must
be recovered from a field on a $3N$-dimensional space, and the manifest world of things in space becomes a
derived appearance. The rival \emph{primitive ontology} programme \cite{Maudlin2007,Allori2013,DGZ1992,Esfeld2014}
insists instead on beables in $\mathbb R^3$ --- particles, a matter density, or flashes --- with $\Psi$
demoted to a dynamical or nomological role. A middle route, \emph{spacetime state realism}, keeps a
description in spacetime by assigning density operators to its regions \cite{WallaceTimpson2010}, yet the
fundamental state remains the global, entangled one and separable objects in $\mathbb R^3$ are not thereby
recovered. Either way the $3N$-dimensional domain is the source of the difficulty: one must swallow it as
fundamental, explain it away, or relocate it.

\emph{Individuality}, from the antisymmetry. The permutation symmetry in \eqref{eq:std} makes the
particles, on the received reading, \emph{non-individuals}: no fact settles which is which, and two
electrons in one atom are numerically distinct yet share all state-dependent and state-independent
properties, in apparent violation of Leibniz's Principle of the Identity of Indiscernibles (PII)
\cite{FrenchKrause2006,FrenchSEP}. Defenders of PII have retreated to \emph{weak discernibility} by
symmetric irreflexive relations, conceding that no distinguishing \emph{monadic} property is available
\cite{Saunders2006,MullerSaunders2008}. Later work has refined the taxonomy of (in)discernibility and
pressed it in both directions \cite{MullerSeevinck2009,CaultonButterfield2012}, but on no variant does an
electron in a many-electron state carry a distinguishing \emph{monadic} property.

Both problems are consequences of the two choices in \eqref{eq:std}: the domain $\mathbb R^{3N}$, and the
antisymmetry. A formulation making neither choice does not incur them.

\section{A real-space ontology, in equations}
\label{sec:po}

Take instead the RealQM formulation \cite{RealQM}
\begin{equation}
\Psi(\mathbf x)=\sum_{i=1}^N\psi_i(\mathbf x),\quad \mathbf x\in\mathbb R^3,\qquad
\operatorname{supp}\psi_i=\overline{\Omega_i},\quad \Omega_i\cap\Omega_j=\varnothing\ (i\neq j),
\label{eq:real}
\end{equation}
where $\psi_i^2$ is the unit charge density of electron $i$ on its own region $\Omega_i\subset\mathbb R^3$
and $\rho=\sum_i\psi_i^2$ is the total electron density. For a nucleus (kernel) of charge $Z$ at the
origin, with attractive potential $V(\mathbf x)=Z/|\mathbf x|$, the ground state minimises the energy
\begin{equation}
E=\sum_i\tfrac12\!\int_{\Omega_i}|\nabla\psi_i|^2
-\int V\rho\,dx
+\tfrac12\!\iint\frac{\rho(x)\rho(y)}{|x-y|}\,dx\,dy,\qquad
\int_{\Omega_i}\psi_i^2\,dx=1,
\label{eq:E}
\end{equation}
over the $\psi_i$ and the free boundaries $\partial\Omega_i$, which meet in Bernoulli (zero-flux) contact
conditions. The Pauli principle enters not as an antisymmetry of a function but as the geometric
requirement that the $\Omega_i$ do not overlap.

RealQM is, at once, a \emph{primitive ontology in three-space}: what exists is the extended charge density
$\rho$ in $\mathbb R^3$, resolved into the $N$ pieces $\psi_i^2$, and there is \emph{no} object on
$\mathbb R^{3N}$ to interpret or to recover three-space from --- \eqref{eq:std}'s domain is simply absent.
Unlike primitive-ontology theories that carry a matter density \emph{alongside} a configuration-space
$\Psi$ that still does the dynamical work \cite{Allori2013}, here the density is the whole matter ontology
and \eqref{eq:E} is its law: there is no second, higher-dimensional object in the theory at all. It is a
rival theory, not a reinterpretation of \eqref{eq:std}, and its warrant is the chemistry and nuclear
structure it computes (\S\ref{sec:comp}).

\section{Individuation by occupation}
\label{sec:ind}

The basic objects of \eqref{eq:real} are extended and disjoint, and this individuates them. Define electron
$i$ as the density on $\Omega_i$; then occupancy of $\Omega_i$ is a \emph{monadic} property that $i$ has
and $j\neq i$ lacks, since $\Omega_i\cap\Omega_j=\varnothing$. Hence
\begin{equation}
i\neq j\ \Longleftrightarrow\ \Omega_i\neq\Omega_j:\qquad
\text{no two electrons share all properties.}
\end{equation}
The Principle of the Identity of Indiscernibles, standardly taken to fail for \eqref{eq:std}, is
\emph{satisfied}. Moreover the electrons are \emph{absolutely} discernible --- each has a distinguishing
monadic property, the region it occupies --- which is stronger than the \emph{weak} discernibility by
relations \cite{MullerSaunders2008} to which PII's defenders retreat under \eqref{eq:std}. Individuality in
the full classical sense, with a haecceity grounded in a spatial fact rather than a primitive thisness, is
recovered.

Extension and individuation are here \emph{the same fact}. The disjointness $\Omega_i\cap\Omega_j=\varnothing$
that keeps each $\psi_i$ a density in $\mathbb R^3$ (the primitive-ontology virtue) is the very disjointness
that discerns the electrons (the individuality virtue). Where \eqref{eq:std} encodes exclusion as
antisymmetry --- a fact about a configuration-space function --- \eqref{eq:real} encodes it as
impenetrability, a fact about \emph{space}: two charge clouds cannot occupy one place. Pauli exclusion and
the identity of indiscernibles, distinct and even opposed in the received picture, are one geometric fact
here.

\section{Computation: where the ontology earns its weight}
\label{sec:comp}

A metaphysics can be clean and still idle. What keeps the present proposal from being one more homemade
picture is that its ontology is also a working computational method, and one that sits closer to the actual
practice of quantum chemistry than the official $3N$-dimensional ontology does.

\paragraph{The dimensional wall.} The amplitude \eqref{eq:std} lives on $\mathbb R^{3N}$, and representing
it there costs exponentially in the number of particles: a modest grid of $m$ points per axis needs
$m^{3N}$ values, so solving \eqref{eq:std} directly is feasible only for very small $N$. Essentially all of
computational quantum chemistry is a response to this wall. Hartree--Fock replaces $\Psi$ by a single
antisymmetrised product and solves in $\mathbb R^3$; post-Hartree--Fock methods add back correlation at
steep cost; and density-functional theory, by the Hohenberg--Kohn theorems \cite{HohenbergKohn1964,
KohnSham1965}, proves that the ground-state energy is a functional of the \emph{three-dimensional} density
$\rho(\mathbf r)$ alone and computes with that object, not with the $3N$-dimensional amplitude. The
overwhelmingly used methods thus already retreat from configuration space to effective three-dimensional
quantities. The $3N$-dimensional object is where the metaphysics is placed but not, for many electrons,
where the computation is done.

\paragraph{RealQM computes with its beables.} In \eqref{eq:real}--\eqref{eq:E} the beables \emph{are} the
objects of computation. One minimises the energy \eqref{eq:E} directly over the charge densities $\psi_i$
and their free boundaries in $\mathbb R^3$; complexity scales with mesh points in three dimensions, not
with the exponential dimension of $\mathbb R^{3N}$. Where density-functional theory makes the
three-dimensional density fundamental as a \emph{calculational} device while leaving the $3N$-dimensional
$\Psi$ as the official ontology, RealQM takes the three-dimensional densities to be fundamental \emph{full
stop}, and reads \eqref{eq:E} as the law they obey. The ontology is not supplied after the calculation, as
an interpretation of a formalism solved elsewhere; it is what one computes with. An ontology that is also
the algorithm is answerable to computation in a way a pure interpretation is not.

\paragraph{The empirical record.} The functional \eqref{eq:E}, with the appropriate kernels, has been
applied across a wide range of systems as a single method, not tuned per case
\cite{RealQM,RealQMatoms,RealQMgallery}:
\begin{itemize}
\item \emph{Atoms.} Ground-state energies computed across the periodic table by direct free-boundary
minimisation, with shell structure emerging from the packing of the $\Omega_i$ rather than from imposed
antisymmetry \cite{RealQMatoms}.
\item \emph{Molecules.} Covalent bonds, molecular geometries, and hydrogen-bonded complexes; for water the
computed dipole moment comes out near the experimental value, and hydrogen-bond energetics are reproduced
\cite{RealQM,RealQMgallery}.
\item \emph{Nuclei, by Coulomb alone.} Read as protons and electrons bound by the same electrostatic energy
\eqref{eq:E} at femtometre scale --- with no strong or weak force --- the model returns nuclear binding
energies of the right size: the $^4$He binding comes out close to experiment, and the ratio of $^4$He to
deuteron binding is reproduced to within a few percent \cite{RealNucleus,QCDBinding}. The same Coulomb
picture describes fusion as a rearrangement of protons and electrons \cite{SolarFusion}.
\end{itemize}
These are not paper derivations but a public library of interactive, browser-runnable simulations
\cite{RealQMgallery} in which the charge densities of atoms, molecules, reactions, and nuclei can be
computed and watched directly, and several of the results are set out in companion papers under review
\cite{RealQMatoms,RealNucleus,QCDBinding,SolarFusion}.

\paragraph{Computing the world is evidence of tracking it.} That a single, parameter-poor functional
reproduces observations across such disparate regimes --- atomic ground-state energies, molecular geometry
and hydrogen bonding, nuclear binding by Coulomb alone --- spanning length scales that differ by some five
orders of magnitude, is not interpretive bookkeeping but hard empirical evidence. A model free to fit one
regime by adjustment is cheap; a model that computes many \emph{independent} regimes at once, under the same
law \eqref{eq:E} and with no per-case tuning, is expensive to fake. On the familiar no-miracles reasoning,
such consilience is a strong sign that the model has latched onto real structure rather than onto an
artefact of representation --- and here that structure just is the ontology of \S\ref{sec:po}: extended
charge in three-space obeying \eqref{eq:E}. Behind this lies a realist principle worth stating plainly:
\emph{what is real is computable}. The ground is the character of physical process itself. Nature runs as an
analog computation --- physical states evolving in real space under physical law --- and an analog process
of this kind can be expected to be reproducible in digital form, by a finite constructive procedure,
provided the mathematical model used expresses the real physics rather than an artefact of representation.
Digital computability is thus the expected signature of a model that captures what the world actually does,
and persistent intractability a warning that a model carries structure the world does not. On the realist
reading nature's analog computation is moreover carried out \emph{in three-dimensional space}, on charge and
field --- not in a $3N$-dimensional configuration space --- so a model formulated in $\mathbb R^3$ shares the
very substrate on which the world computes. That is why RealQM is directly computable there, just where the
$3N$-dimensional amplitude is not: the tractability gap tracks the difference between a model that expresses
the real physical process and one laden with representational overhead. RealQM earns its ontology by
computing what nature computes, in the space where nature computes it; it is on that record, not on
interpretive convenience, that the metaphysics of \S\S\ref{sec:po}--\ref{sec:ind} rests.

\section{The point and the field: a perspective}
\label{sec:sm}

The extension at issue reaches past non-relativistic quantum mechanics. The Standard Model's Lagrangian
carries three gauge forces together with a Higgs and Yukawa sector,
\begin{equation}
\mathcal L_{\rm SM}=-\tfrac14 F^a_{\mu\nu}F^{a\mu\nu}
+\bar\psi\,i\gamma^\mu D_\mu\psi+|D_\mu\Phi|^2-V(\Phi)-\big(y\,\bar\psi_L\Phi\psi_R+\text{h.c.}\big),
\end{equation}
its fundamental carriers \emph{point} excitations of quantum fields. That is a second departure from
\emph{res extensa}: where \eqref{eq:std} moved the object into $\mathbb R^{3N}$, field theory shrank it to
a point in $\mathbb R^3$. The point carries the classical infinity of a zero-extension charge's
self-energy, softened but not cured by renormalisation --- and it is telling that string theory restores
extension (to strings) precisely to tame it. Retaining only the electromagnetic term of $\mathcal L_{\rm
SM}$ and reading it, non-relativistically and electrostatically, as extended densities returns the energy
functional \eqref{eq:E}: RealQM is, in this sense, that one term read as \emph{res extensa} in
$\mathbb R^3$ \cite{RealQMvsSM}. The ontological choice --- the point or the extended thing --- is thus one
question asked at every scale, and RealQM answers it, in its domain, on the side of extension. Nothing here
disputes the Standard Model's quantitative account of sub-nucleonic phenomena; the claim is about which
ontology one reads off the electromagnetic sector at atomic and nuclear scale.

\section{Objections and replies}
\label{sec:obj}

\paragraph{Non-overlap is ad hoc.} \emph{Objection:} imposing $\Omega_i\cap\Omega_j=\varnothing$ is a
stipulation put in to get individuation for free. \emph{Reply:} it is the geometric form of the Pauli
principle, and it is what the free-boundary minimisation of \eqref{eq:E} actually enforces --- the domains
and their boundaries are solved for, not drawn by hand. That the same condition also individuates is the
content of \S\ref{sec:ind}, not an extra assumption: one fact, read twice.

\paragraph{What about interference and entanglement?} \emph{Objection:} the power of \eqref{eq:std} is that
it captures superposition and entanglement, which a sum of disjoint densities seems to lose. \emph{Reply:}
this marks the scope of the claim. RealQM is offered as an ontology and computational model for
\emph{stationary bound states} --- the ground and excited states of atoms, molecules, and nuclei --- where
it is tested \cite{RealQM,RealQMatoms,RealNucleus}. Time-dependent and multi-centre phenomena are treated
by a complex, current-carrying extension of the same real-space formulation \cite{RealQM}; whether that
extension recovers the full range of interference effects is open and is flagged as such (\S\ref{sec:concl}).
The metaphysical thesis here concerns the bound-state regime, and is conditional on empirical adequacy
there.

\paragraph{Is computability a metaphysical virtue at all?} \emph{Objection:} tractability is pragmatic, not
ontological; a theory is not truer for being easier to compute. \emph{Reply:} the point is not ease but
\emph{answerability}. In RealQM the posited beables are the very quantities carried through the calculation
and compared with experiment, so the ontology is directly exposed to empirical test; in the interpretive
programmes the beables are added to a formalism that is solved, and confirmed, in terms of other objects.
An ontology that does the computing is falsifiable in a way an interpretation grafted onto a shared
formalism is not --- and that is an epistemic virtue, not merely a convenience.

\paragraph{It is a fringe, self-published programme.} \emph{Objection:} the physics is not established in
the mainstream literature. \emph{Reply:} granted, and the metaphysical verdict is stated as conditional on
that physics (\S\ref{sec:concl}). But the conditional is not idle: the computational results exist, are
public and reproducible \cite{RealQMgallery}, and are set out in papers under review
\cite{RealQMatoms,RealNucleus,QCDBinding}. The philosophical interest is that \emph{if} a real-space
charge-density model is empirically adequate in its domain --- an open question being pressed --- then it
resolves the arena and individuality problems together, which is a reason to take the model seriously
rather than to dismiss it in advance.

\section{A note on dissolution}
\label{sec:wittgenstein}

There is a deflationary way to hear the foregoing that should be met head on. On the reading defended here
the arena and individuality problems are not so much \emph{solved} as made to \emph{disappear}: once the
fundamental object is extended charge in three-space rather than an amplitude on configuration space, the
question of how to recover ordinary space from $\mathbb R^{3N}$, and the question of what individuates
permutation-symmetric particles, simply do not arise. This is dissolution in Wittgenstein's sense --- the
problems stand revealed as artefacts of a representation mistaken for the world, and clarity makes them
vanish rather than answering them on their own terms \cite{Wittgenstein1921,Wittgenstein1953}. To a
philosopher of physics this can disappoint at first: a substantial and ingenious literature has grown around
puzzles that the realist reformulation quietly removes.

But the disappointment is only the first step. What the dissolution clears away is a set of problems
\emph{internal to a formalism}; what it leaves standing is the harder and more properly physical question of
\emph{what reality is} --- that it should be extended charge in three-dimensional space, evolving under
\eqref{eq:E}, and computable because the world itself computes (\S\ref{sec:comp}). With the classic puzzles
gone, the questions that remain are about the thing itself: the time-dependent, current-carrying regime, the
status of interference and entanglement under a real-space ontology, and why nature's analog computation
takes the form it does. Dissolving the idealist problems is thus not an end but a redirection --- from
puzzles about our representation of matter to the analysis of the matter represented.

\section{Conclusion}
\label{sec:concl}

The arena of the wave function and the individuality of particles are two readings of the two features of
one formula, \eqref{eq:std}: its $3N$-dimensional domain and its antisymmetry. RealQM, the real-space
charge-density formulation \eqref{eq:real}--\eqref{eq:E}, declines both, and so meets both problems with one
move --- a primitive ontology in three-space whose electrons are individuated absolutely, by the regions
they occupy, because $\Omega_i\cap\Omega_j=\varnothing$ does double duty. Extension and individuation are
the same fact, and Pauli exclusion and the identity of indiscernibles coincide. What lifts this above a
merely tidy picture is that the ontology is a working, tested computational method: the standard amplitude
is intractable beyond very small $N$ and practice already retreats to three-dimensional objects, while
RealQM takes those objects as fundamental and computes atoms, molecules, and --- by Coulomb alone --- nuclei
with them \cite{RealQMatoms,RealNucleus,QCDBinding}. The open questions are real: the time-dependent
extension, the full range of interference and entanglement, and the empirical standing of the programme in
the mainstream. But the philosophical payoff is sharp and conditional in a well-defined way: whether such a
model is empirically adequate is for physics to decide; if it is, its metaphysics is, on these two
long-contested points, the cleaner one.

\begin{thebibliography}{99}
\bibitem{Schrodinger1926} E.~Schr\"odinger, \emph{Quantisierung als Eigenwertproblem} (I--IV), Ann.
Phys. \textbf{79}--\textbf{81} (1926); and \emph{An Undulatory Theory of the Mechanics of Atoms and
Molecules}, Phys. Rev. \textbf{28}, 1049 (1926).
\bibitem{Born1926} M.~Born, \emph{Zur Quantenmechanik der Sto\ss vorg\"ange}, Z. Phys. \textbf{37}, 863
(1926).
\bibitem{deRegt1997} H.~W.~de Regt, \emph{Erwin Schr\"odinger, Anschaulichkeit, and the Wave Mechanics},
Studies in History and Philosophy of Modern Physics \textbf{28}, 461 (1997).
\bibitem{Wittgenstein1921} L.~Wittgenstein, \emph{Tractatus Logico-Philosophicus} (1921; Routledge, 1961),
esp. 6.5--6.521.
\bibitem{Wittgenstein1953} L.~Wittgenstein, \emph{Philosophical Investigations} (Blackwell, 1953), esp.
\S\S109--133.
\bibitem{Albert1996} D.~Z.~Albert, \emph{Elementary Quantum Metaphysics}, in J.~Cushing, A.~Fine, and
S.~Goldstein (eds.), \emph{Bohmian Mechanics and Quantum Theory: An Appraisal} (Kluwer, 1996), 277--284.
\bibitem{NeyAlbert2013} A.~Ney and D.~Z.~Albert (eds.), \emph{The Wave Function: Essays on the Metaphysics
of Quantum Mechanics} (Oxford University Press, 2013).
\bibitem{Ney2021} A.~Ney, \emph{The World in the Wave Function} (Oxford University Press, 2021).
\bibitem{Maudlin2007} T.~Maudlin, \emph{The Metaphysics Within Physics} (Oxford University Press, 2007).
\bibitem{Allori2013} V.~Allori, \emph{Primitive Ontology and the Structure of Fundamental Physical
Theories}, in \cite{NeyAlbert2013}, 58--75.
\bibitem{DGZ1992} D.~D\"urr, S.~Goldstein, and N.~Zangh\`i, \emph{Quantum Equilibrium and the Origin of
Absolute Uncertainty}, J. Stat. Phys. \textbf{67}, 843 (1992).
\bibitem{FrenchKrause2006} S.~French and D.~Krause, \emph{Identity in Physics: A Historical, Philosophical,
and Formal Analysis} (Oxford University Press, 2006).
\bibitem{FrenchSEP} S.~French, \emph{Identity and Individuality in Quantum Theory}, Stanford Encyclopedia
of Philosophy (2019).
\bibitem{Saunders2006} S.~Saunders, \emph{Are Quantum Particles Objects?}, Analysis \textbf{66}, 52 (2006).
\bibitem{MullerSaunders2008} F.~A.~Muller and S.~Saunders, \emph{Discerning Fermions}, British Journal for
the Philosophy of Science \textbf{59}, 499 (2008).
\bibitem{MullerSeevinck2009} F.~A.~Muller and M.~Seevinck, \emph{Discerning Elementary Particles},
Philosophy of Science \textbf{76}, 179 (2009).
\bibitem{CaultonButterfield2012} A.~Caulton and J.~Butterfield, \emph{On Kinds of Indiscernibility in Logic
and Metaphysics}, British Journal for the Philosophy of Science \textbf{63}, 27 (2012).
\bibitem{WallaceTimpson2010} D.~Wallace and C.~G.~Timpson, \emph{Quantum Mechanics on Spacetime I:
Spacetime State Realism}, British Journal for the Philosophy of Science \textbf{61}, 697 (2010).
\bibitem{Esfeld2014} M.~Esfeld, \emph{The Primitive Ontology of Quantum Physics: Guidelines for an
Assessment of the Proposals}, Studies in History and Philosophy of Modern Physics \textbf{47}, 99 (2014).
\bibitem{HohenbergKohn1964} P.~Hohenberg and W.~Kohn, \emph{Inhomogeneous Electron Gas}, Phys. Rev.
\textbf{136}, B864 (1964).
\bibitem{KohnSham1965} W.~Kohn and L.~J.~Sham, \emph{Self-Consistent Equations Including Exchange and
Correlation Effects}, Phys. Rev. \textbf{140}, A1133 (1965).
\bibitem{RealQM} [Author], \emph{Quantum Mechanics as Multiphase 3D Continuum Mechanics}, manuscript,
2026; [URL withheld for anonymous review].
\bibitem{RealQMatoms} [Author], \emph{A 3D Multiphase Continuum Computational Model for Atoms},
manuscript, 2026 (under review).
\bibitem{RealNucleus} [Author], \emph{Nuclear Binding as Dual Coulomb Confinement, without a Strong or
Weak Force}, manuscript, 2026 (under review).
\bibitem{QCDBinding} [Author], \emph{The Binding Energy Problem QCD Never Solved --- and a Coulomb Model
That Does}, manuscript, 2026 (under review).
\bibitem{SolarFusion} [Author], \emph{Solar Fusion by Coulomb Alone}, manuscript, 2026.
\bibitem{RealQMvsSM} [Author], \emph{RealQM vs the Standard Model: Extension or Point, One Term or Many},
manuscript, 2026.
\bibitem{RealQMgallery} [Author], \emph{The RealQM Gallery: interactive charge-density simulations of
atoms, molecules, reactions, and nuclei}, [URL withheld for anonymous review]
(source: [URL withheld for anonymous review]).
\end{thebibliography}

\end{document}
